\hypertarget{va.-algemeen}{%
\subsection{Va. Algemeen}\label{va.-algemeen}}

\hypertarget{art.-1}{%
\paragraph{Art. 1}\label{art.-1}}

De Linker, eenieder ander persoon die optreedt als rechter -zowel in
eerste aanleg als in beroep- of de voorzitter van een Tribunaal kan
iedere andere persoon laten verwijderen uit de zitting indien deze
onrust veroorzaken tijdens een hoorzitting.

\hypertarget{vb.-beroep}{%
\subsection{Vb. Beroep}\label{vb.-beroep}}

\hypertarget{art.-2}{%
\paragraph{Art. 2}\label{art.-2}}

In beroep zal een Oepsie poepsie worden berecht door een lid van de
Volksvergadering die voor die specifieke zaak verkozen is door de
Volksvergadering.

\hypertarget{art.-3}{%
\paragraph{Art. 3}\label{art.-3}}

In beroep zal een Da was nie zo goe worden berecht door een lid van de
Volksvergadering die voor die specifieke zaak verkozen is door de
Volksvergadering.

\hypertarget{art.-4}{%
\paragraph{Art. 4}\label{art.-4}}

Er is geen beroep mogelijk tegen een arrest van een Tribunaal.

\hypertarget{art.-5}{%
\paragraph{Art. 5}\label{art.-5}}

Bij een beroep of een voorziening in cassatie zal de Procureur dezelfde
persoon zijn als de persoon die optrad als Procureur in eerste aanleg.

\hypertarget{art.-6}{%
\paragraph{Art. 6}\label{art.-6}}

Tegen een vonnis uitgesproken in beroep kan men niet in beroep gaan.

\hypertarget{vc.-aansprekingen-en-ambtskledij}{%
\subsection{Vc. Aansprekingen en
ambtskledij}\label{vc.-aansprekingen-en-ambtskledij}}

\hypertarget{vc-1.-aansprekingen}{%
\subsubsection{Vc-1. Aansprekingen}\label{vc-1.-aansprekingen}}

\hypertarget{art.-7}{%
\paragraph{Art. 7}\label{art.-7}}

De Linker wordt tijdens een hoorzitting aangesproken met Opperheer.

\hypertarget{art.-8}{%
\paragraph{Art. 8}\label{art.-8}}

Alle leden van een Tribunaal, uitgezonderd de Linker, wordt tijdens een
hoorzitting aangesproken met Frank.

\hypertarget{art.-9}{%
\paragraph{Art. 9}\label{art.-9}}

De Procureur wordt tijdens een hoorzitting aangesproken met Jos.

\hypertarget{art.-10}{%
\paragraph{Art. 10}\label{art.-10}}

De eisende partij in burgelijke zaken of de burgelijke partij in
strafzaken wordt tijdens een hoorzitting aangesproken met Mario.

\hypertarget{art.-11}{%
\paragraph{Art. 11}\label{art.-11}}

De verwerende partij in burgerlijke zaken of de verdediging in
strafzaken wordt tijdens een hoorzitting aangesproken met Luigi.

\hypertarget{vc-2.-ambtskledij}{%
\subsubsection{Vc-2. Ambtskledij}\label{vc-2.-ambtskledij}}

\hypertarget{art.-12}{%
\paragraph{Art. 12}\label{art.-12}}

De ambtskledij wordt gedragen bij elke fysieke hoorzitting door de
Linker, de leden van een Tribunaal, de Procureur, de Griffier en de
advocaten van de aanwezige partijen.

\hypertarget{art.-13}{%
\paragraph{Art. 13}\label{art.-13}}

De basis van de ambtskledij bestaat voor iedereen uit de nationale
klederdracht. Het Hawaii hemd word opengedragen met daaronder ofwel een
wit marcelleke of niks afhankelijk van de omgevingstemperatuur en
voorkeur van de drager.

\hypertarget{art.-14}{%
\paragraph{Art. 14}\label{art.-14}}

De Linker en de leden van een Tribunaal dragen hierbij ook een pruik met
lang haar.

\hypertarget{art.-15}{%
\paragraph{Art. 15}\label{art.-15}}

De Procureur en de advocaten van de aanwezige partijen dragen hierbij
ook een pruik met kort haar.

\hypertarget{art.-16}{%
\paragraph{Art. 16}\label{art.-16}}

De advocaten van de aanwezige partijen dragen hierbij ook een opplaksnor
om hun status als Mario en Luigi aan te duiden.
